\documentclass{article}
\textheight 23.5cm \textwidth 15.8cm
%\leftskip -1cm
\topmargin -1.5cm \oddsidemargin 0.3cm \evensidemargin -0.3cm
%\documentclass[final]{siamltex}

\usepackage[UTF8]{ctex}
\usepackage{verbatim}
\usepackage{fancyhdr}
\usepackage{graphicx}
\usepackage{subcaption}
\usepackage{geometry}
\geometry{a4paper,scale=0.8}

%\pagestyle{fancy} \lhead{FDM Homework Template} \chead{}
%\rhead{\bfseries Yan XU}
%
%\lfoot{} \cfoot{} \rfoot{\thepage}
%\renewcommand{\headrulewidth}{0.4pt}
%\renewcommand{\footrulewidth}{0.4pt}
%
\title{HW1 说明文档}
\author{李佩优}

\begin{document}
\maketitle

\section{引言}

本次作业的任务是对函数$f(x)=\frac{1}{1+x^2}, x \in [-5,5]$ 构造lagrange插值多项式，并分析其$L^{\infty}$范数意义下的误差。

\section{原理}

Lagrange插值多项式的构造方法如下：
用给定插值点$x_1,x_2,...,x_n$，
构造对应基函数$l_i(x)=\prod_{j=1,j\neq i}^n \frac{x-x_j}{x_i-x_j}$；
再线性组合各基函数，构成拉格朗日插值函数$P_L(x)=\sum f(x_i)\cdot l_i(x)$

\section{结果}

\begin{figure}[htb]
\begin{center}
\includegraphics[width=\textwidth]{hw1_btty.png}
\caption{在北太天元中运行N=10时的图像} \label{hw1_btty}
\end{center}
\end{figure}

\begin{figure}[htbp]
    \centering
    % 第一个子图
    \begin{subfigure}[b]{0.45\textwidth}
        \centering
        \includegraphics[width=\textwidth]{result1.png}
        \caption{matlab中运行输出}
        \label{fig:sub1}
    \end{subfigure}
    \hfill % 在子图之间添加间距
    % 第二个子图
    \begin{subfigure}[b]{0.45\textwidth}
        \centering
        \includegraphics[width=\textwidth]{result2.png}
        \caption{北太天元中运行输出}
        \label{fig:sub2}
    \end{subfigure}
    % 整个图集的标题
    % \caption{这是包含两个子图的整体标题}
    \label{fig:total}
\end{figure}

\begin{table}[htb]
\caption{\label{table.label} 误差估计} \centering
\bigskip
\begin{small}
\begin{tabular}{|c|cc|cc|}
  \hline
  % after \\: \hline or \cline{col1-col2} \cline{col3-col4} ...
  n & grid(1)$ L^\infty$ 误差 & grid(2)$ L^\infty$ 误差 \\\hline
5 &   4.327E-01  &  5.559E-01\\
10&   1.916E+00  &  1.089E-01\\
20&   5.827E+01  &  1.533E-02\\
40&   7.869E+04  &  2.739E-04\\\hline
\end{tabular}
\end{small}
\end{table}
\section{讨论}

从结果可以看出，等距节点的误差随着n的增加呈现爆炸式增长，
而切比雪夫节点的误差则随着n的增加呈现指数级下降。
这是因为等距节点在区间两端会出现Runge现象，导致插值多项式在区间两端出现剧烈震荡，从而导致误差增大。
而切比雪夫节点则能够有效地避免Runge现象，使得插值多项式在整个区间内都能够较好地逼近原函数，从而使得误差较小。

\appendix
\section{源代码}


\small \hrule \verbatiminput{hw1.m} \hrule


\end{document}
